\scanpage{107}
\mspage{94}
% archive sheet number: 106

\[
 \bigl(\underset{TC'_{0}}{\underline{u_{!}(TC_{0})}}\bigr)^{\sim}=(u_{!}(TC))^{\sim},
\]
puisque $TC$ est la clôture de $TC_{0}$ pour
les opérations de cochgt. de base, \deleted{\unsure{et}} de
composition transfinie et de facteurs directs,
et que $u_{!}$ commute à ces opérations. Donc
on doit montrer finalement
\[
 TC'\subset (u_{!}(TC))^{\sim}\tag*{(5.8.12 bis)}
\]
(où il n'est plus question de $TC_{0}$).
Mais le th. de factorisation s'applique
au couple
\[
 (u_{!}(TC))^{\sim},\ (u_{!}(TC))_{*}=(u^{*})^{-1}(F)\big), 
\]
puisque $(u_{!}(TC))^{\sim}=(TC'_{0})^{\sim}$, $TC'_{0}$ petit.
Si donc on a $f'\in TC'=C'\cap W'$,
\unsure{factorisant|factorisons} $f'$ en $p'i'$,
\diagram{\begin{tikzcd}[ampersand replacement=\&,column sep=small,row sep=large,execute at end picture={\draw (current bounding box.north west) -- (current bounding box.north east) -- (current bounding box.south east) -- (current bounding box.south west); \node[anchor=east] at ([xshift=-.8em]current bounding box.west) {$(\alpha)$};}]
X' \arrow[rr,dash,"{f'\in C'\cap W'}"] \arrow[dr,"{i'\,\unsure{\in|\notin}}"'] \& \& Y' \arrow[dl,dashed,shift right=1.2ex,"{q'}"']\\
\& X'_{1} \arrow[ur,"{p'}"'] \&
\end{tikzcd}}
avec \deleted{$i'\in (u_{!}(TC))^{\sim}$, \unsure{$p'\in$} $p'\in(u^{*})^{-1}(F)$, \unsure{Noto}}
\[
 i'\in (u_{!}(TC))^{\sim},\qquad p'\in (u^{*})^{-1}(F).\tag*{(5.8.13.)}
\]
Notons qu'on a \deleted{(\unsure{puisque}) à cause de (5.8.9)}
\deleted{\[
 (u_{!}(TC))^{\sim}\subset \deleted{\unsure{C'}}(TC')^{\sim}
\]}
\deleted{et d'autre part je dis qu'on a}
\deleted{\[
 (TC')^{\sim}\subset TC'\deleted{\unsure{\cap W'}}\ (=\unsure{C|C'}\cap W')\tag*{(5.8.14)}
\]}
\deleted{i.e. $TC'$ est Q-clos à \unsure{g.|G.|dr.} En effet, comme \unsure{on a} d'une part $(\underset{C'\cap W'}{TC'})^{\sim}\subset {C'}^{\sim}=C'$, de sorte qu'il suffit de voir que}

